This section reports results from 30{,}000 debates across GPQA and HiddenBench.
Because debate plays a fundamentally different role on each benchmark, every result
is reported per benchmark rather than pooled.
\autoref{sec:results:perf} establishes the accuracy and efficiency baseline for the system itself.
\autoref{sec:results:motifs} measures the four self-organisation motifs against their
nulls and links each to accuracy and efficiency, which addresses RQ1 and RQ2.
\autoref{sec:results:evid} then tests causally whether the opening vote composition
determines the outcome, which addresses RQ3.

\subsection{Accuracy and Efficiency}
\label{sec:results:perf}

Before examining the self-organisation patterns, we report what the system
achieves in aggregate and which design choices shape that outcome.
Accuracy is the fraction of repetitions in which the final plurality vote matches
the ground truth, and efficiency is the mean number of rounds to consensus.
Our reference point for debate gain is the \emph{majority-vote baseline},
which is the plurality of the four agents' round-0 answers before any exchange occurs.

\begin{table}[H]
\centering
\small
\begin{tabular}{lcccc}
\toprule
Dataset & Per-agent & Majority vote (r\,=\,0) & After debate & Oracle (r\,=\,0) \\
\midrule
GPQA        & $38.2\,\%$ & $40.5\,\%$ & $\mathbf{42.2\,\%}$ & $70.6\,\%$ \\
HiddenBench & $14.9\,\%$ & $\phantom{0}7.6\,\%$ & $\mathbf{26.1\,\%}$ & $46.3\,\%$ \\
\bottomrule
\end{tabular}
\caption[Baseline and debate accuracy by dataset]{Accuracy decomposition at round~0
         and after debate, pooled over all configurations with $n = 15{,}000$ per
         dataset. Per-agent is the mean fraction of correct individual answers at
         round~0. Majority vote is the accuracy of the round-0 plurality vote and
         serves as the reference point for all debate-gain calculations. Oracle is
         the fraction of repetitions in which at least one agent is correct at
         round~0.}
\label{tab:perf:baseline}
\end{table}

\autoref{tab:perf:baseline} shows that the two benchmarks operate in qualitatively
different regimes.
On GPQA the round-0 majority vote is already $40.5\,\%$, and debate adds a modest
$+1.7$\,pp to reach $42.2\,\%$. The oracle bound of $70.6\,\%$ shows that the correct
answer is present among the agents in roughly two-thirds of repetitions. The resulting
oracle gap \citep{bertalanic2026costofconsensus} is therefore not a knowledge deficit
but the group's failure to identify and amplify the correct voice.
HiddenBench is different. Here a shared model bias makes majority voting
actively harmful and drags the round-0 plurality well below chance. Debate partly
corrects this and reaches $26.1\,\%$, a gain of $+18.5$\,pp. Repetitions that converge
early lock in the shared wrong answer almost universally, because early consensus
forecloses the corrective exchange that the benchmark design requires, as
\autoref{fig:app:perf:rounds_acc} in Appendix~\ref{app:perf:configs} shows.
In summary, debate refines a mostly-correct start on GPQA. On HiddenBench it partially
corrects a systematically wrong one.

With the outcome picture in place, the remaining question is what design choices such
as memory window and topology contribute. \autoref{fig:perf:memory_window} plots
accuracy and rounds across all configurations.

\begin{figure}[H]
  \centering
  \includegraphics[width=0.92\linewidth]{plots/general/fig3_memory_window.png}
  \caption[Accuracy and rounds by memory window]{Accuracy in the top row and mean
           rounds to consensus in the bottom row as a function of memory window
           $W \in \{1,2,5\}$, for GPQA in the left column and HiddenBench in the right
           column. Each line is one topology, and the gap between the lines in the
           rounds panels is the topology efficiency effect. The y-axes do not start
           at zero, and the accuracy differences are not significant.}
  \label{fig:perf:memory_window}
\end{figure}

Neither structural axis has a significant effect on accuracy. Topology makes
no detectable difference on either benchmark. Memory window shows a
consistent pattern, as $W=2$ yields the highest accuracy for both datasets and
both topologies. However, the differences are small and do not survive
multiple-comparison correction, so we read $W=2$ as the best setting in our data
rather than as an established effect.

Both axes do affect efficiency. The star topology needs more rounds to reach
consensus on both benchmarks, and especially on HiddenBench, with $p < 0.001$ in
both cases. The reason is that every spoke must route its messages through a single
hub. The star partly compensates by sending half as many messages per round, so each
round costs fewer tokens. Under the star topology, a larger memory window also speeds
up convergence on HiddenBench. Exact values for all configurations are given in
Appendix~\ref{app:perf:configs}.


% -----------------------------------------------------------------------------
\subsection{Self-Organization Motifs}
\label{sec:results:motifs}

Each motif is reported with a three-step template. Detection gives the prevalence
against a null, drivers cover topology, $W$, and the initial vote composition, and
outcomes relate the motif to accuracy and rounds to consensus.

\subsubsection{Self-Reinforcement}
\label{sec:results:sr}

Self-reinforcement asks whether an agent that keeps its vote grows steadily more
confident in it. We read it from the slope $\beta_{\mathrm{seg}}$ of confidence
against round number within each \emph{vote-stable segment}, which is a run of at
least three consecutive rounds on the same vote. We summarise a cell by the prevalence
$p_{\mathrm{SR}}$, the fraction of segments with a positive slope, against a
random-walk null of $0.5$. The full construction is given in \autoref{sec:method:sr}.
Pooled across both benchmarks and all configurations, the segment filter leaves
$133{,}696$ segments. Three findings follow.

\begin{enumerate}
\renewcommand{\labelenumi}{(\roman{enumi})}
  \item \textbf{It is universal.} $p_{\mathrm{SR}}$ never approaches its null. It
        ranges from $0.733$ to $0.881$ across the 12 combinations of dataset,
        topology, and $W$, and all 50 tasks lie individually above $0.5$ in every
        combination. Per-configuration values are given in
        Appendix~\ref{app:sr:supp}. Confidence climbs fastest where the debate has
        least to resolve, in repetitions that never change their votes and in
        unanimous starts. The initial vote composition drives the slope far more
        than topology or memory window do. Agents grow more certain in what they
        already held, not in what the debate revealed.

  \item \textbf{It drives convergence.} That rising confidence is what makes the
        debate settle. Tasks with steeper confidence slopes reach consensus in fewer
        rounds in all four combinations of dataset and topology. The task-level
        Spearman correlation lies between $\rho = -0.82$ and $\rho = -0.73$, as
        \autoref{fig:sr:efficiency} shows.

  \item \textbf{It is blind to correctness.} Faster convergence is not more accurate
        convergence. Whether self-reinforcement helps depends entirely on \emph{what}
        it amplifies. When the correct option is already among the opening votes,
        stronger reinforcement pushes the group toward it. When no agent starts
        correct, the same force only locks in a wrong answer sooner.
        Appendix~\ref{app:sr:supp} gives the corresponding accuracy analysis.
\end{enumerate}

\begin{figure}[H]
  \centering
  \includegraphics[width=\linewidth]{plots/sr/fig4_sr_efficiency.png}
  \caption[Self-reinforcement slope versus rounds to consensus]{Mean segment slope
           $\bar\beta$ against mean rounds to consensus at the task level. Each panel
           is one combination of dataset and topology, each point is one task, and
           colours denote the memory window. The dashed line is the pooled OLS fit.}
  \label{fig:sr:efficiency}
\end{figure}

Self-reinforcement converges the debate, but it converges on whatever the opening votes
already favour rather than on the correct answer. This matches the position lock-in
that \citet{oh2026entrenchment} identify as a root cause of belief fixation, that
\citet{estornell2024multillm} derive from response homogeneity, and that
\citet{cau2026agreementdrift} describe as a selective amplifier of locally consistent
change. It also mirrors reports that debate confidence rises without new evidence
\citep{prasad2025certain, bertalanic2026costofconsensus}. The seeded experiment in
\autoref{sec:results:evid} confirms this dependence on the initial composition
causally.

% -----------------------------------------------------------------------------
\subsubsection{Multistability}
\label{sec:results:multistability}

Multistability asks whether a task admits more than one stable consensus outcome, and
whether the opening vote distribution decides which one the debate reaches. We read
it per task-configuration cell from two signatures. The first is $N_{\mathrm{eff}}$,
the inverse-Simpson diversity of the endpoint attractors, which counts how many stable
outcomes coexist across repetitions. The second is Cramér's $V$, which measures how
strongly the initial votes select among these outcomes. A cell is \emph{monostable}
when one attractor dominates with $N_{\mathrm{eff}} < 1.5$, \emph{multistable} when
several coexist and the initial votes predict the endpoint, and \emph{stochastic}
when several coexist but the endpoint does not follow from the start. The full
construction is given in \autoref{sec:method:multistability}. Only repetitions that
settle on a unanimous endpoint are included. Three findings follow.

\begin{enumerate}
\renewcommand{\labelenumi}{(\roman{enumi})}
  \item \textbf{Most tasks are multistable, and the structure is set by the question.}
        Across the 600 task-configuration cells, $76\,\%$ show more than one
        coexisting attractor with $N_{\mathrm{eff}} > 1.5$. This rate is conditioned
        on the selected task set. Neither topology nor memory window shifts either
        metric, and per-task $N_{\mathrm{eff}}$ is highly consistent across the six
        configurations. The structure is therefore a property of the question.
        Multiple absorbing states are theoretically expected here. Naming-game
        dynamics admit one equilibrium per available answer option
        \citep{baronchelli2018consensus, ashery2025conventions}, so a four-option
        task can in principle land on any of four attractors. Prior LLM debate
        studies could not observe this because they report a single run or
        run-level means per task \citep{cau2026agreementdrift,
        mehdizadeh2026topologymemory}. Our repeated-runs design makes the full
        attractor landscape visible.

  \item \textbf{Multistability helps exactly when the default answer is wrong.} Across
        tasks, more coexisting attractors go with \emph{higher} accuracy on
        HiddenBench ($\rho = +0.44$) but slightly \emph{lower}, non-significant
        accuracy on GPQA ($\rho = -0.16$). The opposition follows from what the
        dominant attractor is. On GPQA, where agents usually start near the correct
        answer, a single attractor is already the right one, and extra basins only
        add wrong-answer traps. On HiddenBench, where the majority attractor is the
        shared wrong answer, competing basins give the correct minority a foothold
        from which to win.

  \item \textbf{Fewer attractors produce stronger self-reinforcement.} Monostable cells
        carry the steepest confidence ramps, and $N_{\mathrm{eff}}$ anticorrelates
        with the mean self-reinforcement slope at the task level ($\rho = -0.75$ on
        GPQA and $\rho = -0.44$ on HiddenBench). When one attractor dominates, agents
        lock in early and confidence climbs unchecked. When several coexist, the
        reinforcement signal is split across diverging trajectories.
        Self-reinforcement is the within-basin amplifier, and multistability counts
        the basins.
\end{enumerate}

Multistability is thus the structural counterpart of the amplifier. Each task carries a
small set of stable outcomes fixed by the question itself. Which of those outcomes the
group reaches depends on where it starts, as the causal analysis in
\autoref{sec:results:evid} will show. A diversity of attractors therefore helps only when
the single basin the group would otherwise collapse into is the wrong one. The per-dataset breakdown of monostable, multistable, and stochastic cells is given in Appendix~\ref{app:multistab}.


% -----------------------------------------------------------------------------
\subsubsection{Dominance and Effective Topology}
\label{sec:results:dominance}
\label{sec:results:eff_topology}

Dominance measures how unequally influence is shared among the agents. We compute it
from the confidence-weighted conversion-credit graph as the Herfindahl index $D$. A
value of $0$ means that every agent converts equally many peers, and a value of $1$
means that a single agent causes every opinion change. The full construction is given
in \autoref{sec:method:dominance}.
One structural caveat limits what we can test. In the star topology, a leaf can only
ever be converted by the centre. Its raw $D$ is therefore inflated by the wiring alone,
and its surrogate null is degenerate for most repetitions. We report raw $D$ for both
topologies, but we test for \emph{emergent} concentration only in the fully-connected
topology, where all agents are structurally equal.
We also describe the \emph{shape} of the emerging hierarchy. For this, we map the
conversion-credit graph of each fully-connected run to one of four structural types.
In a \emph{flat} run, influence is nearly uniform. A \emph{hub} run has one moderately
dominant agent, and a \emph{two-hub} run has two agents of comparable influence. In a
\emph{star} run, a single agent is responsible for at least $75\,\%$ of all
conversions. \autoref{fig:efftopo:examples} shows an example of each type. Four
findings follow.

\begin{figure}[H]
\centering
\includegraphics[width=\linewidth]{plots/eff_topology/fig1_topology_examples.pdf}
\caption[Effective topology types in FC runs]{%
  Examples of the four effective topologies observed in fully-connected runs with
  $N{=}4$. Arrow width is proportional to vote-flip influence weight. The orange
  node marks the hub, which is the most influential agent. Green borders mark
  agents that held the correct answer at round~0.}
\label{fig:efftopo:examples}
\end{figure}

\begin{enumerate}
\renewcommand{\labelenumi}{(\roman{enumi})}
  \item \textbf{A hierarchy emerges without a hub.} Every agent in a fully-connected
        run is wired identically, yet influence is almost never uniform. Conversion
        credit concentrates above the topology-keyed null in 49 of 50 GPQA tasks and
        in all 50 HiddenBench tasks. Flat runs are rare. By far the most common type
        is two-hub, in which two agents of comparable influence act as competing
        faction leaders. This type typically arises from an initial vote split of
        two against two. A clean star is rare, as \autoref{tab:efftopo:types} shows.

  \item \textbf{A leader helps only when the debate selects it.} When a star forms
        in a fully-connected debate, the group has chosen a hub that often started
        with the correct answer. On GPQA this holds for $54.5\,\%$ of star hubs and
        on HiddenBench for $85.5\,\%$. A random agent is correct at round~0 only
        $38.2\,\%$ and $14.9\,\%$ of the time, as \autoref{tab:perf:baseline} shows.
        The group then almost always adopts the hub's answer. If the hub started
        correct, about nine in ten runs end correct. If it started wrong, fewer than
        one in five do. This is why effective stars reach the highest accuracy of all
        types, with a clear margin on HiddenBench, as \autoref{tab:efftopo:types}
        shows. The picture reverses when the leader is imposed by the wiring. In the
        star topology on HiddenBench, the correlation between dominance and accuracy
        is negative at $\rho = -0.29$. The fixed centre is no more likely to be
        correct than any other agent, so it spreads the shared initial bias instead
        of correcting it.

  \item \textbf{Diverse initialisation is the lever that produces stars.} Whether a
        star forms depends strongly on how diverse the opening votes are. Stars
        almost never form when all agents start on the same option. When all four
        agents start on different options, roughly one in five runs becomes a star,
        as \autoref{fig:efftopo:design} shows. Enforced perspective diversity has
        already been proposed as a remedy for belief lock-in
        \citep{oh2026entrenchment}, and it also follows from the theoretical analysis
        of debate homogeneity \citep{estornell2024multillm}. Our contribution is the
        mechanism behind it. Diversity raises the chance that the debate elects an
        informed hub.

  \item \textbf{Concentration costs rounds.} Whatever its effect on accuracy,
        dominance consistently slows convergence. Within each task, configurations
        with a higher mean $D$ need more rounds to reach consensus. The per-task
        Spearman correlation is $\rho = +0.55$ on GPQA and $\rho = +0.73$ on
        HiddenBench. A dominant agent has to press its position over many rounds
        before the other agents give in.
\end{enumerate}

\begin{table}[H]
\centering
\small
\begin{tabular}{lcccc}
\toprule
 & \multicolumn{2}{c}{Share of FC runs} & \multicolumn{2}{c}{Accuracy} \\
\cmidrule(lr){2-3}\cmidrule(lr){4-5}
Effective topology & GPQA & HiddenBench & GPQA & HiddenBench \\
\midrule
Flat         & $5.4\,\%$  & $6.8\,\%$  & $0.394$ & $0.242$ \\
Hub          & $11.0\,\%$ & $18.0\,\%$ & $0.393$ & $0.377$ \\
Two-hub      & $58.7\,\%$ & $59.2\,\%$ & $0.418$ & $0.247$ \\
Star         & $2.3\,\%$  & $5.4\,\%$  & $\mathbf{0.485}$ & $\mathbf{0.562}$ \\
Unclassified & $22.6\,\%$ & $10.6\,\%$ & $0.450$ & $0.029$ \\
\addlinespace
All FC runs  &            &            & $0.423$ & $0.264$ \\
\bottomrule
\end{tabular}
\caption[Effective topology types, prevalence and accuracy]{Prevalence and accuracy
         of the four effective topology types in fully-connected runs, with
         $7{,}500$ runs per dataset. Unclassified runs contain no vote flips and
         therefore have no conversion-credit graph.}
\label{tab:efftopo:types}
\end{table}

\begin{figure}[H]
\centering
\includegraphics[width=\linewidth]{plots/eff_topology/fig3_design_guideline.pdf}
\caption[Initial vote diversity and star topology emergence]{%
  Probability of an effective star topology as a function of initial vote diversity.
  $n_\text{unique}$ is the number of distinct votes held at round~0. The
  maximum-diversity bin contains $144$ GPQA runs but only $6$ HiddenBench runs.}
\label{fig:efftopo:design}
\end{figure}

Together, dominance and effective topology show the amplifier from two complementary
angles. The debate reliably elects an influence hierarchy that the wiring never
specified, and the group's answer follows whichever agent wins that election. Accuracy
is high when the elected hub started correct, and it suffers when a fixed hub
entrenches a shared error. This is the structural counterpart to the causal test in
\autoref{sec:results:evid}, which shows the same dependence on the opening votes by
direct manipulation. It also mirrors reports that nominally symmetric LLM populations
develop structural asymmetry on their own, both through preferential-attachment links
\citep{mehdizadeh2025homophily} and through locally clustered agents coordinating better than high-betweenness bridges
\citep{mehdizadeh2026topologymemory}. Per-cell concentration values across all
configurations are in Appendix~\ref{app:dom:supp}.


% -----------------------------------------------------------------------------
\subsubsection{Limit Cycles}
\label{sec:results:lc}

Limit cycles are the rarest motif in this system, and the only one that is a failure
rather than a variant of convergence. A limit cycle is a debate that never settles and
oscillates between vote configurations instead of reaching a fixed point. We detect
them with categorical recurrence quantification. A sequence is flagged when its
determinism exceeds a shuffle-null floor and its dominant period lies in a genuine
oscillatory band with $2 \leq \hat{P} \leq \lfloor L/2 \rfloor$. The full construction
is given in \autoref{sec:method:lc}. The cooperative design suppresses oscillation
almost entirely. We test at two levels, the vote sequence of each individual agent
and the vote composition of the whole group per repetition. At the agent level,
$99.5\,\%$ of sequences converge to a fixed point, and at the system level
$98.7\,\%$ do. These converged sequences never enter the test. Three findings follow.

\begin{enumerate}
\renewcommand{\labelenumi}{(\roman{enumi})}
  \item \textbf{They are rare but genuine.} Within the small non-converged remainder,
        47 agent-level and 32 system-level sequences are flagged, both far above the
        shuffle-null floor. On inspection, $70\,\%$ of the flagged agent-level cases
        are clean period-2 orbits whose recurrence matrices show the characteristic
        checkerboard pattern. The remaining cases are near-convergent artefacts.

  \item \textbf{They occur only in the star topology.} Almost every flagged case sits
        in a star configuration, 45 of 47 at the agent level and 26 of 32 at the
        system level. The few flagged fully-connected cases all turn out to be stable
        revisits or near-convergent artefacts on inspection. Flagged cases also
        concentrate on HiddenBench, with 34 of 47 at the agent level.
        \autoref{fig:lc:timelines} shows the two clearest examples.

  \item \textbf{The mechanism is hub-driven majority inversion, purely an efficiency
        failure.} Every limit cycle runs to the hard round cap without ever reaching
        consensus, so early stopping never fires. Inspection suggests a consistent,
        though uncertified, causal chain. Agents defer to the peer majority they
        observe. After each flip, the star hub rebroadcasts that majority for the
        opposite option. With two independently defensible lines of evidence and a
        hub that starts in the minority, each flip recreates the conditions that
        caused it. This also explains why the fully-connected topology shows no
        genuine cycles. Without a single node that aggregates all peer signals, the
        inversion loop cannot form. In practice, repetitions that reach the cap
        should therefore be flagged as unresolved rather than reported as a
        decision. Failure to recognise termination conditions is a catalogued failure
        mode of deployed multi-agent systems \citep{cemri2025mast}, and our result
        names a specific dynamical route to it.
\end{enumerate}

\begin{figure}[H]
  \centering
  \includegraphics[width=\linewidth]{plots/lc/lc_timelines_new.png}
  \caption[Vote-sequence timelines for two limit cycles]{Vote-sequence timelines
         for two limit cycles in the star topology. Each cell shows the vote held
         by an agent in a given round, and colour encodes the option. Both panels
         show the same hub-driven majority inversion, in which three agents
         alternate in phase and one agent alternates in anti-phase. In (a), from
         GPQA with $W{=}2$ and ground truth B, the cycle starts almost immediately.
         In (b), from HiddenBench with $W{=}5$ and ground truth A, the group first
         agrees on B and only falls into the cycle from round~5 onward. Neither
         repetition reaches consensus, and both run to $T = 15$.}
  \label{fig:lc:timelines}
\end{figure}

Limit cycles are therefore the failure mode of the same convergence machinery seen
elsewhere. When a single aggregator keeps reconstructing the opposite majority,
self-reinforcement never gets to lock the group in, and the debate oscillates instead
of settling. That this happens only under the star, and never when every agent sees
every other, fits the observation that unanimous states are otherwise effectively
absorbing \citep{li2024sparsetopology}. The full detection funnel and
representative recurrence matrices are given in Appendix~\ref{app:lc:supp}.

% -----------------------------------------------------------------------------
\paragraph{Motif summary.}
\autoref{tab:motif:summary} collects the four motifs side by side. Three of them
point back to the opening vote composition. Self-reinforcement amplifies it,
multistability lets it select the basin, and the formation of an elected hub
depends on its diversity. Limit cycles are the exception, since they arise from
the star wiring rather than from the opening votes. \autoref{sec:results:evid}
then manipulates the opening vote composition directly.

\begin{table}[H]
\centering
\small
\begin{tabularx}{\textwidth}{
  >{\raggedright\arraybackslash}p{2.6cm}
  >{\raggedright\arraybackslash}X
  >{\raggedright\arraybackslash}X
  >{\raggedright\arraybackslash}X
  >{\raggedright\arraybackslash}p{2.2cm}}
\toprule
Motif & Detection & Main driver & Accuracy & Rounds \\
\midrule
Self-reinforcement
  & In every cell and task
  & Opening votes
  & Helps only if the correct answer is present at the start
  & Fewer \\
\addlinespace
Multistability
  & In most cells
  & The task itself
  & Higher on HiddenBench, no significant link on GPQA
  & Not tested \\
\addlinespace
Dominance and effective topology
  & Above the null in fully-connected runs
  & Diversity of the opening votes
  & Helps if the hub is elected, hurts if it is imposed
  & More \\
\addlinespace
Limit cycles
  & Rare, star topology only
  & Star hub
  & None, the debate stays unresolved
  & Always hits the cap \\
\bottomrule
\end{tabularx}
\caption[Summary of self-organisation motif findings]{Summary of the four motif
         findings. Detection refers to prevalence against the respective null,
         and accuracy and rounds describe the direction of the association.}
\label{tab:motif:summary}
\end{table}


% -----------------------------------------------------------------------------
\subsection{Causal Test of the Initial-Condition Effect}
\label{sec:results:evid}

The motif analyses point to the opening vote composition as the factor that decides
which attractor the debate reaches. This association is confounded, because whether
the ground truth is present at $t=0$ co-varies with task difficulty, agent confidence,
and initial vote entropy. Three experiments address this in sequence. A within-task
observational analysis at $R=1{,}000$ establishes the pattern. A seeded dose-response
fixes the opening composition directly and tests the causal interpretation. An
adversarial probe then stress-tests the robustness of the amplifier.

\paragraph{Ground-truth presence in the initial vote distribution.}
For the observational analysis, three GPQA questions, q84, q125, and q144, were each
run for $R = 1{,}000$ repetitions at $W = 2$ in the fully-connected topology with
$N = 4$. They were selected for mid-range accuracy and high initial-vote entropy, so
that their outcomes are genuinely uncertain. The dominant finding is striking. Whether
the ground-truth option appears at all in the initial vote distribution is by far the
strongest predictor of final accuracy, consistently across all three tasks, as
\autoref{fig:highrep:init_votes_split} and \autoref{tab:highrep:init_votes} show.

\begin{figure}[H]
  \centering
  \includegraphics[width=0.9\linewidth]{plots/highrep/fig1_init_votes_split.png}
  \caption[Accuracy by ground-truth presence at initialisation]{System accuracy
           conditioned on whether the ground-truth option is present in the initial
           vote distribution at $t=0$. Left bars show repetitions in which at least
           one agent voted correctly before debate, and right bars show repetitions
           in which no agent did. The dashed line marks the mean per-agent accuracy
           at $t=0$. Error bars are Wilson 95\,\% confidence intervals
           \citep{wilson1927ci}. The data cover three GPQA tasks with $R=1{,}000$
           each, $W=2$, and the fully-connected topology.}
  \label{fig:highrep:init_votes_split}
\end{figure}

\begin{table}[H]
\centering
\small
\begin{tabular}{lcccc}
\toprule
Task & Overall & GT absent & GT present & GT is plurality \\
     & acc.    & (acc., $n$) & (acc., $n$) & (acc., $n$) \\
\midrule
q84  & 0.269 & 0.004~(257) & 0.361~(743) & 0.620~(213) \\
q125 & 0.273 & 0.031~(293) & 0.373~(707) & 0.552~(250) \\
q144 & 0.509 & 0.069~(144) & 0.583~(856) & 0.717~(403) \\
\bottomrule
\end{tabular}
\caption[Accuracy by ground-truth presence per task]{Accuracy by ground-truth presence
         in the initial vote distribution, with $R=1{,}000$ per task. GT absent and
         GT present indicate whether the correct option appears in the vote profile
         at $t=0$. GT is plurality is the subset of present repetitions in which the
         correct option holds the initial plurality.}
\label{tab:highrep:init_votes}
\end{table}

When the correct answer is absent from all initial votes, the system almost never
reaches it. Accuracy is $0.4\,\%$, $3.1\,\%$, and $6.9\,\%$ for q84, q125, and
q144, with $\chi^2 > 120$ and $p < 10^{-27}$ in all cases.
The rare correct answers in the absent condition are consistent with stochastic
noise rather than genuine recovery.
Within the present condition, accuracy is roughly twice as high when the correct
option also holds the initial plurality. For q84, for example, accuracy is $0.620$
when the correct option leads the opening votes and $0.257$ when it is present but
does not lead.
These findings suggest that the system operates primarily as a vote amplifier.
Debate selectively strengthens options already represented in the initial
distribution and overwhelmingly reinforces whichever option holds the initial
plurality. The capacity to surface a correct answer that no agent endorsed at the
outset is negligible under the present design.

Neither the length of a debate nor its confidence trajectory predicts whether it is
right. Longer debates are no more accurate, and confidence rises monotonically across
rounds whether the final answer is correct or not. Appendix~\ref{app:highrep} gives
the full feature-by-outcome correlations. What the outcome does track is the
opening vote composition, which the next experiment manipulates directly.

% -----------------------------------------------------------------------------
\paragraph{Seeded dose-response.}
\label{sec:results:seeded}
To isolate the causal contribution of the opening composition, we fix it directly.
We use $50$ GPQA tasks in the same configuration with $W=2$, the fully-connected
topology, $N=4$, and $R=50$ repetitions each. For every repetition, we build the
initial four-agent vote profile by sampling from a pool of genuine, independent
single-agent responses under a fixed number of correct seeds $k \in \{0,1,2,3\}$.
Exactly $k$ of the four agents begin on the correct option, and the rest begin on
real wrong answers. Only the initial composition is manipulated. Every vote is a real
response, and the debate that follows is unmodified. Sweeping $k$ turns the earlier
binary contrast into a dose-response curve.

\begin{figure}[H]
  \centering
  \includegraphics[width=\linewidth]{plots/highrep/fig4_seeded_causal.png}
  \caption[Causal dose-response of seeding correct agents]{Accuracy as a function of
           the number of correct agents seeded into the initial vote profile,
           $k \in \{0,1,2,3\}$, on $50$ GPQA tasks with $R=50$ each, $W=2$, and the
           fully-connected topology. Panel (a) shows per-task accuracy in grey, one
           line per task, and the mean over tasks in orange. Panel (b) shows pooled
           accuracy in blue against the
           accuracy that a pure plurality vote at $t=0$ would already achieve, shown
           as a dashed grey line. Debate exceeds this amplifier baseline for a
           correct minority or tie but falls below it for a correct supermajority.}
  \label{fig:highrep:seeded_causal}
\end{figure}

Three findings follow.

\begin{enumerate}
\renewcommand{\labelenumi}{(\roman{enumi})}
  \item \textbf{Accuracy rises monotonically with the number of correct seeds.}
        Pooled over $2{,}500$ repetitions per condition, accuracy climbs from
        $1.8\,\%$ at $k=0$ to $26.2\,\%$, $63.7\,\%$, and $86.3\,\%$ at $k=1,2,3$,
        as \autoref{fig:highrep:seeded_causal} shows. This is not merely a pooled
        average. Adding one correct seed raises per-task accuracy in $50$ of $50$
        tasks for each of the first two steps and in $47$ of $50$ for the last.
        Causally, the composition of the opening votes is the dominant lever on the
        outcome.

  \item \textbf{Debate is a lossy amplifier, not a reasoner.} Comparing each
        condition against the accuracy that a pure plurality vote at $t=0$ would
        already achieve isolates what the debate itself contributes. This baseline
        is the dashed line in \autoref{fig:highrep:seeded_causal}. Where the correct
        answer starts behind, debate helps. A lone correct minority at $k=1$, where
        plurality accuracy is $0\,\%$, is carried to the right answer $26\,\%$ of the
        time. An even split of two against two at $k=2$, with a baseline of
        $38\,\%$, resolves correctly $64\,\%$ of the time. Where the correct answer
        starts ahead, debate \emph{erodes} it. A correct supermajority of three at
        $k=3$, which plurality voting would preserve with certainty, is talked out of
        the right answer in roughly one repetition in seven, reaching $86\,\%$ rather
        than $100\,\%$. From a standing start, debate generates almost nothing. It
        recovers the answer de novo in only $1.8\,\%$ of $k=0$ repetitions, confined
        to $21$ of $50$ tasks. Debate amplifies the initial signal while blurring it
        toward a coin flip, rather than reasoning its way to the truth.

  \item \textbf{Outcomes are all-or-nothing and decided early.} Every repetition in
        every condition ends fully unanimous. The correct seeds either convert the
        whole group or capitulate entirely, and no repetition ends with the ground
        truth as a surviving minority. When the correct side wins, it wins fast and
        reaches a plurality within the first one to two rounds, with a mean first
        round of $1.3$. Persuasion is front-loaded. The opening exchange, not
        protracted deliberation, decides the outcome.
\end{enumerate}

% -----------------------------------------------------------------------------
\paragraph{Adversarial probe.}
The amplifier does not care whether it amplifies a correct or a wrong answer, and
this becomes most visible when one agent acts as an adversary. We re-run the $k=2$
condition with $R=25$, but now one of the two agents without the correct answer is
told to keep arguing for a wrong option throughout the debate. Accuracy drops from
$63.7\,\%$ to $10.0\,\%$, and it drops in every single one of the $50$ tasks. The group
now does worse than a single correct agent without any adversary, which still reaches
$26.2\,\%$ at $k=1$. This happens even though the adversary starts with two correct
agents against it. Most of these runs still reach unanimity, and when they do, the
group almost always ends on the adversary's wrong answer.
The reason is simple. An agent that never changes its mind looks convinced, and the
debate rewards conviction. The same mechanism that sometimes lets a correct minority
win now works for the adversary, so the group follows the most persistent voice
rather than the most correct one. This matches the tipping-point dynamics that
\citet{centola2018tipping} report for human conventions, where a committed minority of
about $25\,\%$ is enough to overturn the majority. Our adversary is exactly such a
committed minority of one in four agents. Similar effects have been reported for LLMs.
\citet{agarwal2025cwpor} show that a forcefully argued falsehood can override a truthful answer in an LLM judge's decision, and \citet{yao2025sycophancy} describe how agents give in to
each other too early, a problem that prompt-based fixes now try to address
\citep{pitre2025consensagent}.

Taken together, the three experiments upgrade the amplifier reading from correlation
to intervention. The multi-agent system is a \emph{lossy amplifier of its initial vote
distribution}. It reliably strengthens whatever the opening votes favour, more so as
that signal grows, yet it neither manufactures a correct answer that no agent held
nor protects a correct one that a persistent dissenter attacks. The composition of
the opening votes, not the length or intensity of the debate, is the dominant lever
on whether collective reasoning succeeds. This intervention was run on GPQA only. The
observational ground-truth-presence result holds on both benchmarks, as
\autoref{sec:results:sr} shows, and the extension of the causal claim to HiddenBench
is discussed in \autoref{sec:discussion:limitations}.